\documentclass[12pt, a4paper]{article}

\usepackage[utf8]{inputenc}
\usepackage[T1]{fontenc}
\usepackage{lmodern}

\usepackage[spanish, es-nodecimaldot, es-tabla]{babel}


\usepackage[top=2.25cm, bottom=2cm, left=3cm, right=2.5cm]{geometry}

\usepackage[style=ieee, backend=biber, sorting=none, doi=true]{biblatex}
\addbibresource{referencias.bib}

\usepackage{setspace}
\onehalfspacing
\setlength{\parindent}{1.25cm}
\setlength{\parskip}{2pt}

\usepackage{titlesec}
\titleformat{\section}{\large\bfseries}{\thesection.}{0.5em}{}
\titleformat{\subsection}{\normalsize\bfseries}{\thesubsection.}{0.5em}{}
\titleformat{\subsubsection}{\normalsize\itshape}{\thesubsubsection.}{0.5em}{}
\titlespacing{\section}{0pt}{14pt}{6pt}
\titlespacing{\subsection}{0pt}{10pt}{4pt}
\titlespacing{\subsubsection}{0pt}{8pt}{3pt}


\usepackage{booktabs}
\usepackage{array}
\usepackage{float}
\usepackage{caption}
\captionsetup{font=small, labelfont=bf, skip=6pt}

\usepackage[hidelinks, breaklinks=true]{hyperref}
\usepackage{url}
\urlstyle{tt}

\usepackage{fancyhdr}
\fancyhf{}
\fancyhead[L]{Aplicaciones Distribuidas \textperiodcentered{} UTEQ}
\fancyhead[R]{\small Cristhian Daniel Pacheco Cárdenas}
\fancyfoot[C]{\thepage}
\renewcommand{\headrulewidth}{0.4pt}
\setlength{\headheight}{14pt}

\usepackage[framemethod=default]{mdframed}
\usepackage{xcolor}
\definecolor{azulUTEQ}{RGB}{13, 42, 93}
\definecolor{grisClaro}{RGB}{245, 245, 245}
\definecolor{azulClaro}{RGB}{232, 240, 254}

\usepackage{enumitem}
\setlist[itemize]{noitemsep, topsep=3pt, leftmargin=1.4em}
\setlist[enumerate]{noitemsep, topsep=3pt, leftmargin=1.6em}

\begin{document}

\begin{titlepage}
\centering

\vspace*{0.5cm}

{\large\bfseries Universidad Técnica Estatal de Quevedo}\\[6pt]
{\normalsize Facultad de Ciencias de la Computación y Diseño Digital}\\[4pt]
{\normalsize Carrera de Ingeniería en Software}

\vspace{1cm}
\rule{\linewidth}{1pt}\\[10pt]

{\Large\bfseries
  Análisis comparativo entre la Arquitectura Orientada
  a Servicios (SOA) y la Arquitectura de Microservicios:\\[4pt]
  Justificación de la decisión de diseño en el Sistema\\[4pt]
  de Gestión de Solicitudes de Soporte Técnico de Internet
}

\vspace{2pt}
\rule{\linewidth}{1pt}

\vspace{1.2cm}

{\normalsize
\textbf{Nombre:} Cristhian Daniel Pacheco Cárdenas\\[4pt]
  \textbf{Asignatura:} Aplicaciones Distribuidas\\[4pt]
  \textbf{Tipo de actividad:} Trabajo autónomo individual sumativo\\[4pt]
  \textbf{Período académico:} Regular 2026--2027 SPA\\[4pt]
  \textbf{Semestre:} Séptimo Semestre\\[4pt]
  \textbf{Docente:} Gleiston Cicerón Guerrero Ulloa
}

\vspace{4.5cm}

{\small Quevedo, Ecuador \textperiodcentered{} 13 de junio de 2026}

\end{titlepage}


\pagestyle{fancy}


\begin{abstract}
\noindent
El presente informe desarrolla un análisis comparativo entre la Arquitectura
Orientada a Servicios (SOA) y la arquitectura de microservicios, con el objetivo
de fundamentar la decisión de diseño adoptada en el Proyecto Fin de Curso (PFC):
el Sistema de Gestión de Solicitudes de Soporte Técnico de Internet. El sistema
evolucionó desde una arquitectura en tres capas esencialmente monolítica---
hacia un ecosistema de cinco microservicios desacoplados, comunicados mediante
REST síncrono y Apache Kafka asíncrono. Se examinan los principios fundacionales
de cada estilo, sus atributos de calidad bajo el marco ISO/IEC~25010 y las
condiciones concretas que determinan la aplicabilidad de uno u otro. Los hallazgos
demuestran que la adopción de microservicios en lugar de SOA responde a tres
restricciones verificables del proyecto: fronteras de dominio bien delimitadas
por diseño guiado por el dominio (DDD), disponibilidad de infraestructura de
contenedores y ausencia de sistemas heterogéneos preexistentes que justifiquen
un ESB. Se concluye que los microservicios con mensajería orientada a eventos
satisfacen los requisitos de escalabilidad, tolerancia a fallos y desplegabilidad
independiente exigidos por el dominio de soporte técnico, mientras que SOA
habría introducido complejidad sin resolver ningún problema real del sistema.

\medskip
\noindent\textbf{Palabras clave:} arquitectura distribuida, SOA, microservicios,
ESB, API Gateway, Apache Kafka, DDD, bounded context, Circuit Breaker, Saga,
atributos de calidad, soporte técnico.
\end{abstract}

\newpage
\tableofcontents
\newpage

\section{Introducción}
\label{sec:intro}

El diseño de sistemas distribuidos plantea una de las decisiones más determinantes
en ingeniería de software: elegir el estilo arquitectónico sobre el que se
construirá el sistema. Esa decisión no es fácilmente reversible una vez que el
código existe y el equipo se ha organizado alrededor de un modelo de trabajo.
Por eso, tomarla con criterios técnicos verificables y no con base en
tendencias de mercado es una responsabilidad central del arquitecto de software.

El Sistema de Gestión de Solicitudes de Soporte Técnico de Internet denominado
en adelante \textit{el sistema} o PFC nació como una aplicación web monolítica
en tres capas. Esta arquitectura inicial tiene ventajas reconocibles: es simple de
razonar, fácil de depurar y suficiente para un equipo pequeño en fase de exploración
del dominio. Sin embargo, a medida que el sistema incorpora un componente de
inteligencia artificial para la clasificación automática de incidencias, notificaciones
multicanal y reportes analíticos, la arquitectura monolítica se convierte en una
limitación estructural~\cite{blinowski2022}.

La pregunta que motiva este informe no es si migrar a una arquitectura más
distribuida esa decisión ya está tomada en el PFC sino cuál de los dos
estilos candidatos responde mejor a las necesidades concretas del sistema: la
Arquitectura Orientada a Servicios (SOA) o la arquitectura de microservicios.
Ambas proponen distribuir la lógica en unidades de negocio autónomas; ambas buscan
bajo acoplamiento y alta cohesión. Pero difieren profundamente en cómo logran esos
objetivos y qué precio operativo cobran por ellos~\cite{shadija2017}.

SOA integra servicios a través de un Bus de Servicios Empresariales (ESB) que
centraliza el enrutamiento, la transformación de mensajes y la orquestación de
procesos. Los microservicios, en cambio, distribuyen esa responsabilidad: cada
servicio decide cómo comunicarse con los demás, sin intermediario central, usando
protocolos ligeros y mensajería asíncrona~\cite{liu2020}. La diferencia parece
sutil en el papel, pero sus implicaciones operativas son radicales en la práctica.

\subsection{Objetivo general}

Comparar de manera crítica los estilos arquitectónicos SOA y microservicios para
fundamentar la decisión de diseño adoptada en el Sistema de Gestión de Solicitudes
de Soporte Técnico de Internet.

\subsection{Objetivos específicos}

\begin{enumerate}
  \item Caracterizar los principios fundacionales de SOA y microservicios a partir
        de literatura académica verificada.
  \item Contrastar ambas arquitecturas en criterios de escalabilidad, despliegue,
        gestión de datos, comunicación y tolerancia a fallos.
        \newpage
  \item Proyectar los hallazgos del análisis sobre el contexto real del PFC,
        identificando cuál estilo resuelve mejor sus restricciones verificables.
  \item Reconocer las limitaciones de la arquitectura adoptada y proponer medidas
        de mitigación concretas.
\end{enumerate}

\subsection{Alcance}

El análisis cubre los estilos SOA y microservicios en su dimensión de diseño
lógico. No evalúa infraestructura de red, hardware ni proveedores de nube
específicos. La proyección al PFC toma como referencia la Entrega~2 del sistema,
que establece el diseño de arquitectura de microservicios con cinco servicios
independientes, comunicados mediante REST y Apache Kafka.


\section{Marco teórico}
\label{sec:marco}


\subsection{Arquitectura de software de un sistema distribuido}

La arquitectura de software de un sistema distribuido es el conjunto de
\textbf{decisiones estructurales} que determinan qué componentes existen, qué
responsabilidad tiene cada uno y cómo se comunican, con independencia del hardware
donde se ejecuten. Esta concepción distingue dos vistas: la \textit{vista física}
nodos, redes y protocolos de transporte y la \textit{vista lógica} omponentes, módulos, servicios, contratos e interfaces.

El presente análisis trabaja sobre la vista lógica porque es la que determina las
propiedades de mantenibilidad, escalabilidad y desplegabilidad. Como señalan
Shadija, Rezai y Hill~\cite{shadija2017}, la arquitectura lógica de un sistema
distribuido condiciona directamente sus atributos de calidad bajo el marco
ISO/IEC~25010: escalabilidad, disponibilidad, tolerancia a fallos y desplegabilidad.

\subsection{El monolito en capas como punto de partida}

Antes de comparar SOA y microservicios, es necesario situar el punto de origen del
sistema: una arquitectura en tres capas compuesta por una capa de presentación
la interfaz web, una capa de lógica de negocio validaciones, flujos de
asignación y gestión de roles y una capa de datos el repositorio relacional
con toda la información del sistema. Las tres capas se despliegan como una sola
unidad ejecutable sobre una base de datos compartida.

Las ventajas de este punto de partida son reales: simplicidad inicial, transacciones
locales con consistencia inmediata, depuración de extremo a extremo sencilla y
latencia mínima entre capas porque las llamadas ocurren en memoria~\cite{blinowski2022}.
Sin embargo, a medida que el sistema crece incorporando inteligencia artificial,
reportes analíticos y notificaciones multicanal el monolito acumula tres problemas
estructurales: el escalado es en bloque, cualquier cambio requiere redesplegar toda
la aplicación y la base de datos compartida impide la evolución independiente de
los módulos.

\subsection{Arquitectura Orientada a Servicios}

SOA es un estilo en el que las capacidades de negocio se exponen como
\textbf{servicios reutilizables} con contratos formales típicamente WSDL y
mensajes SOAP integrados a través de un Bus de Servicios Empresariales (ESB).
El ESB actúa como intermediario central que enruta, transforma y orquesta mensajes
entre los distintos sistemas de la organización.

Sus principios fundamentales, documentados por Shadija et al.~\cite{shadija2017},
incluyen: acoplamiento débil entre servicios, contrato de servicio estandarizado,
abstracción de la implementación, reutilización y autonomía relativa. En la
práctica, SOA fue adoptado masivamente en entornos empresariales durante la primera
década del siglo~XXI para resolver el problema de integración entre sistemas
heterogéneos ERPs, CRMs y aplicaciones legadas sin necesidad de reescribirlos.

El riesgo histórico de SOA reside precisamente en el ESB: al concentrar lógica de
enrutamiento, transformación de mensajes y orquestación de procesos, el bus se
convierte en un componente de alta complejidad interna. Cuando el ESB falla o se
satura, la comunicación de toda la organización se interrumpe de forma simultánea.
Liu et al.~\cite{liu2020} identifican este punto único de fallo como la razón
principal por la que SOA fue desplazado por los microservicios en contextos de alta
disponibilidad.

\subsection{Arquitectura de microservicios}

Los microservicios son, según Shadija et al.~\cite{shadija2017}, una forma
particular de SOA aplicada al nivel de la aplicación en lugar del nivel empresarial.
La diferencia fundamental no está en los principios ambos buscan servicios con
contratos y bajo acoplamiento sino en la granularidad, el mecanismo de integración
y la cultura de equipo requerida.

Liu et al.~\cite{liu2020} definen la arquitectura de microservicios como un
paradigma que descompone una aplicación en \textbf{unidades funcionales pequeñas
e independientes}, cada una responsable de una capacidad de negocio concreta,
ejecutándose en su propio proceso y comunicándose mediante mecanismos ligeros.
Sus propiedades clave son: modelado alrededor de una capacidad de negocio no de
una capa técnica, despliegue independiente, datos descentralizados con cada
servicio como propietario exclusivo de su base de datos, y tolerancia a fallos
mediante aislamiento de procesos.

Blinowski et al.~\cite{blinowski2022} advierten que los microservicios eliminan el
ESB central distribuyendo la inteligencia de comunicación en los propios servicios,
lo que reduce el acoplamiento estructural al precio de incrementar la complejidad
operativa: los equipos deben gestionar observabilidad distribuida, descubrimiento
de servicios, resiliencia de red y consistencia eventual. Esta advertencia es
directamente relevante para el PFC: adoptar microservicios sin infraestructura
de contenedores equivale a crear un \textit{monolito distribuido} lo peor de
ambos mundos.

\subsection{La Arquitectura Orientada a Eventos como patrón complementario}

La Arquitectura Orientada a Eventos (EDA) actúa como complemento natural de los
microservicios. El productor publica un evento en un intermediario de mensajes
(\textit{broker}) como Apache Kafka y continúa trabajando sin esperar respuesta.
Los consumidores se suscriben y reaccionan al evento a su propio ritmo, produciendo
un desacoplamiento temporal entre servicios.

Rahmatulloh et al.~\cite{rahmatulloh2022} demuestran empíricamente que integrar
EDA dentro de una arquitectura de microservicios mejora el tiempo de respuesta en
un 19{,}18\,\% y reduce la tasa de error en un 34{,}40\,\% frente a la
comunicación HTTP síncrona, si bien incrementa el uso de CPU en un 8{,}52\,\%.
Este hallazgo es directamente aplicable al PFC, donde Kafka conecta de forma
asíncrona el servicio de tickets con el de inteligencia artificial y el de
notificaciones, permitiendo que el cliente reciba el identificador del ticket de
forma inmediata sin esperar a que la IA complete la clasificación.

El estudio de caso de Laigner et al.~\cite{laigner2020} migración real de un
sistema monolítico de Big Data a microservicios con EDA confirma que este patrón
facilita el mantenimiento y el aislamiento de fallos en sistemas con múltiples
actores y flujos de trabajo complejos, que es exactamente el escenario del PFC.

% ============================================================
\section{Análisis comparativo: SOA vs.\ microservicios}
\label{sec:analisis}
% ============================================================

\subsection{Criterios de comparación}

El análisis se estructura en torno a seis dimensiones técnicas extraídas de la
literatura revisada y alineadas con los atributos de calidad ISO/IEC~25010:
unidad de despliegue, escalabilidad, gestión de datos, modelo de comunicación, tolerancia a fallos y complejidad operativa.

\subsection{Unidad de despliegue}

En SOA, la unidad de despliegue es el servicio empresarial, un componente de
tamaño mediano o grande que depende del ESB para ser invocado. Actualizar un
servicio implica coordinar su nueva versión con el bus y verificar compatibilidad
con todos sus consumidores registrados, lo que frena la velocidad de entrega.

En microservicios, cada unidad de despliegue es un proceso independiente. Esto
permite publicar cambios en el servicio de notificaciones sin detener ni modificar
el servicio de tickets ni el de autenticación. Laigner et al.~\cite{laigner2020}
confirman esta ventaja en un estudio de caso real: la migración a microservicios
con EDA facilitó el mantenimiento y el aislamiento de fallos de manera notable,
aunque la diversidad tecnológica incorporada fue identificada como una desventaja
operativa concreta.

\subsection{Escalabilidad}

SOA permite escalar servicios individuales, pero el ESB actúa como limitante: si
el bus se satura, la capacidad de los servicios independientes queda subutilizada.
Los microservicios permiten lo que Blinowski et al.~\cite{blinowski2022} denominan
\textit{escalado fino}: solo el servicio saturado se replica, sin necesidad de
escalar el sistema completo.

No obstante, los mismos autores formulan una advertencia crítica: para sistemas
que no manejan miles de usuarios concurrentes, los beneficios del escalado en
microservicios no han sido investigados con rigor suficiente y la evidencia
disponible es inconsistente~\cite{blinowski2022}. Adoptar microservicios en un
sistema de tráfico moderado puede resultar en mayor complejidad operativa sin
ningún beneficio real de rendimiento, advertencia que el equipo del PFC debe
tener presente durante la fase de implementación.

\subsection{Gestión de datos}

Esta dimensión representa la diferencia más profunda entre ambos estilos. En SOA,
la práctica habitual consiste en compartir una base de datos centralizada entre
múltiples servicios, lo que simplifica las consultas cruzadas mediante sentencias
\texttt{JOIN} pero genera acoplamiento de datos que impide la evolución independiente.

En microservicios, el patrón \textit{Database per Service} es un principio
fundacional: cada servicio es propietario exclusivo de sus datos y ningún otro
puede consultarlos directamente. Los \texttt{JOIN} entre servicios deben
reemplazarse por composición de llamadas REST o eventos. Para gestionar
transacciones distribuidas se emplean los patrones CQRS Segregación de
Responsabilidad entre Comandos y Consultas y Saga secuencia de transacciones
locales con compensaciones ante fallo.

\subsection{Modelo de comunicación}

SOA utiliza predominantemente protocolos pesados: SOAP sobre HTTP con mensajes XML
y contratos WSDL. La transformación y el enrutamiento de mensajes ocurren dentro
del ESB, lo que introduce latencia adicional y rigidez contractual.

Los microservicios favorecen protocolos ligeros: REST sobre HTTP para APIs públicas,
gRPC para comunicación interna de alta frecuencia, y mensajería asíncrona para
eventos y notificaciones. Liu et al.~\cite{liu2020} consolidan la regla práctica
de la industria: REST hacia afuera, gRPC o mensajería hacia adentro, y eventos
para todo aquello que pueda tolerar latencia.
\newpage
\subsection{Tolerancia a fallos}

En SOA, el ESB constituye un punto único de fallo de alto impacto: si el bus cae,
todos los servicios pierden su canal de comunicación simultáneamente. En
microservicios, la tolerancia a fallos se construye mediante patrones específicos:
\textit{circuit breaker} cortacircuito que deja de enviar peticiones a un
servicio degradado, reintentos con retroceso exponencial y tiempos de espera
explícitos. El fallo de un servicio se manifiesta como degradación controlada del
sistema en lugar de un colapso general.

\subsection{Complejidad operativa}

SOA presenta complejidad centralizada en el ESB: configuración de rutas, políticas
de transformación y gobierno de servicios. Los microservicios presentan complejidad
alta y distribuida: contenedores, observabilidad, descubrimiento de servicios y
CI/CD. Liu et al.~\cite{liu2020} identifican la depuración entre servicios, la
latencia acumulada y la consistencia de datos distribuidos como los desafíos más
críticos en producción.

\subsection{Síntesis comparativa}

La Tabla~\ref{tab:comparativa} consolida los resultados del análisis.

\begin{table}[H]
\centering
\caption{Comparativa técnica entre SOA y microservicios}
\label{tab:comparativa}
\begin{tabular}{>{\bfseries}p{3.0cm} p{5.0cm} p{5.0cm}}
\toprule
\textbf{Criterio} & \textbf{SOA} & \textbf{Microservicios} \\
\midrule
Unidad de despliegue
  & Servicio grande coordinado con el ESB
  & Proceso pequeño e independiente \\
\addlinespace
Escalabilidad
  & Por servicio, limitada por el ESB
  & Granular: solo el servicio saturado \\
\addlinespace
Gestión de datos
  & Base de datos frecuentemente compartida
  & Una base por servicio (\textit{Database per Service}) \\
\addlinespace
Comunicación
  & SOAP/XML enrutado por el ESB
  & REST, gRPC, mensajería asíncrona \\
\addlinespace
Tolerancia a fallos
  & ESB como punto único de fallo
  & Aislamiento por servicio (\textit{circuit breaker}) \\
\addlinespace
Complejidad operativa
  & Alta, centralizada en el ESB
  & Alta, distribuida: contenedores, CI/CD, observabilidad \\
\addlinespace
Contexto adecuado
  & Integración de sistemas heterogéneos empresariales preexistentes
  & Equipos autónomos, dominio maduro, despliegue frecuente \\
\bottomrule
\end{tabular}
\end{table}

\subsection{Análisis crítico}

La Tabla~\ref{tab:comparativa} revela que SOA y microservicios son respuestas a
problemas distintos en escalas distintas. SOA resuelve la \textit{interoperabilidad
empresarial}: cómo hacer que sistemas heterogéneos escritos décadas antes intercambien
información sin reescribirlos. Los microservicios resuelven la \textit{velocidad de
entrega y escalabilidad autónoma}: cómo permitir que equipos independientes
desplieguen y escalen sus componentes sin coordinación central.

Un error conceptual frecuente es evaluar los microservicios como una versión
mejorada de SOA. Shadija et al.~\cite{shadija2017} aclaran que los microservicios
son SOA aplicado al nivel de la aplicación. La diferencia no está en los principios
sino en la granularidad, el mecanismo de integración y la cultura de equipo: SOA
centraliza la inteligencia en el bus, los microservicios la distribuyen en los
propios servicios.

La limitación más crítica de los microservicios, señalada con rigor por Blinowski
et al.~\cite{blinowski2022}, es que para sistemas de escala moderada la complejidad
operativa puede superar los beneficios. Un equipo que implementa microservicios sin
contenedores ni observabilidad no obtiene despliegue independiente: obtiene un
sistema distribuido frágil y costoso de mantener. Este riesgo es concreto en
proyectos académicos de equipos de cuatro personas y debe gestionarse con disciplina.

SOA, por su parte, tiene una limitación estructural igualmente seria: el ESB
concentra inteligencia en una capa que no pertenece a ningún dominio de negocio,
y cuando falla, toda la organización lo siente al mismo tiempo. Esa fragilidad fue,
históricamente, la razón que motivó el surgimiento de los microservicios.

\section{Aplicación al Proyecto Fin de Curso}
\label{sec:pfc}


\subsection{Descripción del sistema}

El Sistema de Gestión de Solicitudes de Soporte Técnico de Internet es una
plataforma web orientada a formalizar el proceso de atención de incidencias
relacionadas con la conectividad dentro de entornos organizacionales. El sistema
opera mediante un esquema de acceso restringido por roles: el \textit{cliente}
registra solicitudes y realiza seguimiento; el \textit{técnico} atiende los casos
asignados; el \textit{administrador} supervisa el sistema y asigna técnicos; el
\textit{superusuario} gestiona cuentas y perfiles.

El dominio central es el ciclo de vida completo de un ticket: desde que el cliente
reporta una incidencia pérdida de conexión, degradación de velocidad, fallo de
hardware de red hasta la resolución confirmada y el cierre formal, pasando por
la asignación a técnicos, el registro de trámites, las notificaciones de estado y
el análisis inteligente de patrones de falla.
\newpage
El proyecto es desarrollado por un equipo de cuatro personas con roles definidos:
arquitecto responsable del diseño distribuido, líder de desarrollo responsable de
la implementación de los servicios críticos, responsable de calidad a cargo de las
pruebas de integración y responsable de documentación.

\subsection{Evolución arquitectónica: del monolito a los microservicios}

La versión inicial del sistema se construyó sobre una arquitectura en tres capas:
capa de presentación, capa de lógica de negocio y capa de datos. Esta arquitectura,
apropiada para la fase de exploración del dominio, presentó limitaciones concretas
a medida que el sistema incorporó capacidades adicionales:

\begin{itemize}
  \item El componente de inteligencia artificial para clasificación automática de
        tickets requiere un proceso independiente con dependencias Python,
        incompatibles con el entorno Java del sistema principal.
  \item El servicio de notificaciones multicanal debe poder desplegarse y escalarse
        de forma independiente, ya que su latencia no debe bloquear la creación
        de solicitudes por parte del cliente.
  \item El módulo de reportes analíticos necesita un modelo de datos desnormalizado
        para consultas complejas, incompatible con la base de datos operacional
        compartida.
\end{itemize}

Estas tres necesidades concretas no una tendencia tecnológica justificaron
la migración a microservicios en la Entrega~2 del PFC.

\subsection{Identificación de bounded contexts mediante DDD}

La descomposición en microservicios siguió los principios del Diseño Guiado por
el Dominio (DDD). Se identificaron cinco \textit{bounded contexts} diferenciados
por su responsabilidad:

\begin{itemize}
  \item \textbf{BC-1: Autenticación y Autorización.} Gestiona la identidad de
        todos los actores, emite y valida tokens JWT y controla permisos basados
        en roles (RBAC). Actúa como proveedor de identidad para todos los demás.
  \item \textbf{BC-2: Gestión de Tickets.} Centraliza el ciclo de vida completo
        de una solicitud de soporte. Es el núcleo del dominio (\textit{core domain}).
  \item \textbf{BC-3: Notificaciones.} Consume eventos del bus de mensajes y
        distribuye notificaciones según el canal configurado.
  \item \textbf{BC-4: Inteligencia Artificial.} Clasifica automáticamente las
        incidencias, estima prioridad y sugiere técnicos basándose en historial.
  \item \textbf{BC-5: Reportes y Análisis.} Construye y expone métricas operativas:
        tiempo promedio de resolución, cumplimiento de SLA y volumen de incidencias.
\end{itemize}

\subsection{Catálogo de microservicios adoptados}

La Tabla~\ref{tab:microservicios} describe el catálogo completo de microservicios
del sistema resultante de la migración.

\begin{table}[H]
\centering
\caption{Microservicios del Sistema de Gestión de Soporte Técnico}
\label{tab:microservicios}
\begin{tabular}{>{\ttfamily}p{4.5cm} p{1.5cm} p{3.25cm} p{5.5cm}}
\toprule
\textbf{Servicio} & \textbf{Puerto} & \textbf{Base de datos} & \textbf{Responsabilidad principal} \\
\midrule
auth-service         & 8001 & PostgreSQL & Registro, login, JWT y RBAC \\
ticket-service       & 8002 & PostgreSQL & CRUD de tickets, estados, asignaciones y SLA \\
notification-service & 8003 & MongoDB    & Consumo de eventos Kafka, despacho multicanal \\
ai-service           & 8004 & MongoDB    & Clasificación NLP, priorización y sugerencias \\
report-service       & 8005 & PostgreSQL & Modelo CQRS de lectura, métricas y exportación \\
\bottomrule
\end{tabular}
\end{table}

El API Gateway (Spring Cloud Gateway, puerto~8080) actúa como punto único de
entrada para todos los clientes externos: enruta peticiones, verifica tokens JWT
delegando al \texttt{auth-service}, aplica límites de tasa y centraliza el
registro de accesos. Apache Kafka opera como bus de mensajes asíncrono para los
eventos de dominio entre servicios.

\subsection{Patrones de microservicios aplicados}

El sistema implementa tres patrones que responden a desafíos concretos del dominio:

\textbf{CQRS:} El \texttt{ticket-service} recibe un alto volumen de escrituras
pero el \newline {report-service} necesita consultas complejas con agregaciones que
penalizarían la base de datos operacional. La solución separa el modelo de escritura
del modelo de lectura: el ticket-service actualiza su propia base de datos;
el report-service mantiene una base de datos desnormalizada que se actualiza
escuchando los eventos Kafka del \texttt{ticket-service}.

\textbf{Saga por coreografía.} La creación de un ticket involucra múltiples
servicios en cadena. El ticket-service crea el registro en estado
\texttt{PENDIENTE\_CLASIFICACION} y publica \texttt{ticket.created}; el
\texttt{ai-service} clasifica la incidencia y publica \texttt{ticket.classified};
el \texttt{ticket-service} actualiza categoría y prioridad y publica
\texttt{ticket.ready}; el \newline {notification-service} envía la confirmación al
cliente. Si la clasificación falla, el \texttt{ai-service} publica un evento de
compensación y el ticket queda en estado de error, sin afectar a los demás
servicios. 

\newpage
Este flujo implementa exactamente el desacoplamiento temporal que
Rahmatulloh et al.~\cite{rahmatulloh2022} demuestran que mejora el rendimiento y
reduce la tasa de errores.

\textbf{Circuit Breaker:} Si el \texttt{ticket-service} se degrada, el API Gateway
acumularía peticiones en espera hasta agotar su pool de hilos, dejando inoperables
incluso las rutas de autenticación. El Circuit Breaker implementado con Resilience4j
abre el circuito al superar el 50\,\% de tasa de fallo en las últimas diez
peticiones y devuelve inmediatamente una respuesta controlada sin propagar el
timeout al cliente mientras los demás microservicios continúan operando.

\subsection{Por qué no SOA: justificación técnica}

El rechazo de SOA como alternativa se fundamenta en tres razones técnicas
verificables derivadas del análisis de la Sección~\ref{sec:analisis}:

\textbf{Primera:} el sistema se construye desde cero con tecnología seleccionada
por el equipo. No existe ningún sistema legado heterogéneo que requiera un ESB
para integrarse. Introducir un ESB en este contexto equivale a crear artificialmente
el problema que SOA fue diseñado para resolver.

\textbf{Segunda:} un ESB centralizado contradice el requisito de disponibilidad
del sistema de soporte. Si el bus falla, la comunicación de todos los servicios
se interrumpe simultáneamente. La arquitectura de microservicios con Kafka y
Circuit Breaker garantiza que el fallo de un servicio incluso el de inteligencia
artificial se degrada de forma controlada sin colapsar el sistema completo.

\textbf{Tercera:} los cinco bounded contexts identificados mediante DDD tienen
fronteras suficientemente claras como para justificar servicios autónomos. No
existe coordinación constante entre ellos que requiera un orquestador central.

\subsection{Limitaciones reconocidas y medidas de mitigación}

El equipo reconoce tres limitaciones técnicas concretas:

\begin{enumerate}
  \item \textbf{Consistencia eventual.} Las operaciones que involucran más de un
        servicio no pueden garantizar consistencia inmediata. El patrón Saga mitiga
        este problema, pero introduce lógica de compensación que debe diseñarse y
        probarse con cuidado. Laigner et al.~\cite{laigner2020} confirman que este
        costo fue uno de los más significativos en su migración real. El equipo
        debe documentar los estados transitorios aceptables y los tiempos máximos
        de propagación de cada evento.
\newpage
  \item \textbf{Complejidad de depuración distribuida.} Los errores que atraviesan
        múltiples servicios son más difíciles de rastrear que en un monolito. Como
        mitigación, el equipo debe incorporar identificadores de correlación en los
        encabezados HTTP y en los metadatos de cada evento Kafka, siguiendo la
        recomendación de Liu et al.~\cite{liu2020}, de modo que los logs de los
        cinco servicios puedan correlacionarse durante la depuración.

  \item \textbf{Carga operativa por diversidad tecnológica.} El sistema integra
        Spring Boot (Java), Node.js y FastAPI (Python), con bases de datos
        PostgreSQL y MongoDB. Esta heterogeneidad es una ventaja arquitectónica
        cada servicio usa la tecnología más adecuada pero una carga real
        para un equipo de cuatro personas. Laigner et al.~\cite{laigner2020}
        documentaron exactamente este trade-off en su caso de estudio.
\end{enumerate}

\section{Conclusiones}
\label{sec:conclusiones}

El análisis comparativo entre SOA y microservicios confirma que ninguno de los
dos estilos es universalmente superior: son respuestas a problemas distintos en
contextos distintos. SOA resuelve la interoperabilidad entre sistemas heterogéneos
preexistentes mediante contratos formales y un bus de orquestación; los
microservicios resuelven la velocidad de entrega y la escalabilidad granular
mediante servicios autónomos y comunicación ligera.

Para el Sistema de Gestión de Solicitudes de Soporte Técnico de Internet, la
elección de microservicios es la correcta por tres razones verificables: el
dominio tiene cinco fronteras de negocio limpias identificadas mediante DDD, la
infraestructura de contenedores está disponible con Docker y Docker Compose, y no
existe ningún sistema legado que requiera un ESB para integrarse. SOA habría
introducido complejidad de configuración sin resolver ningún problema real del PFC.

Sin embargo, el equipo debe asumir conscientemente los costos de la decisión:
consistencia eventual en operaciones distribuidas, complejidad de depuración entre
servicios y carga operativa de mantener múltiples pilas tecnológicas. Estos costos
no son argumentos en contra de la arquitectura adoptada; son restricciones de diseño
que deben gestionarse desde las primeras iteraciones con patrones específicos Saga,
identificadores de correlación, builds multi-stage y no descubrir en producción.

La enseñanza más importante del análisis, respaldada por Blinowski
et al.~\cite{blinowski2022}, es que los beneficios de los microservicios no son
automáticos. Solo se materializan cuando el equipo tiene la disciplina operativa
necesaria y el dominio tiene las fronteras claras que permiten que la autonomía de
los servicios sea real. En el PFC, ambas condiciones están presentes. En sistemas
donde no lo están, el monolito modular sigue siendo la decisión más honesta.


\section*{Repositorio en GitHub}

El código en LaTex de este documento se encuentra disponible en:
\newline
\url{https://github.com/CristhianP03/Proyecto-Grupal-App-Distribuidas-EquipoACC/edit/main/PFC%20Entrega%202/Arquitectura%20de%20Microservicios%2C%20API%20REST%20y%20Contenerizaci%C3%B3n%20con%20Docker.TeX}

\section*{Rol de la Inteligencia Artificial}
Durante la elaboración de este informe se utilizó una herramienta de
inteligencia artificial como apoyo en la redacción. Su participación se limitó
a correcciones menores de texto y estructura, como mejorar la redacción de
algunos párrafos, reorganizar el orden de ciertas explicaciones y unificar
el estilo del documento.

\clearpage

\printbibliography[title={Referencias bibliográficas}]

\end{document}
